% PUV Pipeline — Beamer Slide Deck
% Last updated: April 5, 2026
% Metropolis theme

\documentclass[10pt, aspectratio=169]{beamer}
\usetheme{metropolis}

\usepackage{booktabs}
\usepackage{amsmath}
\usepackage{siunitx}
\usepackage{graphicx}
\usepackage{xcolor}
\usepackage{appendixnumberbeamer}
% \usepackage[draft]{todonotes}  % uncomment for \missingfigure support

% ------------------------------------------------------------------
% Metadata
% ------------------------------------------------------------------
\title{PUV Nearshore Wave Measurement Pipeline}
\subtitle{Processing, Validation, and Spectral Bias Investigation}
\author{Holden Bole}
\institute{Scripps Institution of Oceanography}
\date{April 2026}

% Footer with last-updated stamp and frame numbers
\setbeamertemplate{frame footer}{Last updated: April 5, 2026}
\setbeamertemplate{frame numbering}[fraction]

% Shorthand for figure path
\graphicspath{{../outputs/validation/}}

% ------------------------------------------------------------------
\begin{document}

\maketitle

\begin{frame}{Outline}
  \tableofcontents
\end{frame}

% ==================================================================
\section{Overview}
% ==================================================================

% ------------------------------------------------------------------
\begin{frame}{Motivation}
  \begin{itemize}
    \item Nearshore wave transformation governs beach morphology,
          sediment transport, and coastal hazards
    \item Accurate in-situ wave measurements at \SIrange{5}{15}{\meter}
          depth are needed to validate spectral wave models and drive
          transport calculations
    \item PUV (pressure--velocity) measurements from bottom-mounted ADVs
          provide directional spectra without surface moorings
    \item Goal: build a robust, reproducible processing pipeline from
          raw Nortek Vector data to validated wave statistics
  \end{itemize}
\end{frame}

% ------------------------------------------------------------------
\begin{frame}{Instrument Network}
  \begin{columns}[T]
    \begin{column}{0.55\textwidth}
      \begin{itemize}
        \item \textbf{20 deployments}, \textbf{40 instruments}
        \item 5 sites along San Diego coast:
          \begin{itemize}
            \item Torrey Pines --- MOP 580 \& 586
            \item SIO Pier --- MOP 511
            \item Solana Beach --- MOP 654
            \item Los Pe\~nasquitos Lagoon --- MOP 591
          \end{itemize}
        \item Depths: \SIrange{5}{15}{\meter}
        \item Deployment period: 2023--2026
        \item Nortek Vector ADVs, \SI{2}{\hertz} continuous
      \end{itemize}
    \end{column}
    \begin{column}{0.42\textwidth}
      % TODO: Insert map figure when available
      % \includegraphics[width=\textwidth]{site_map.png}
      \centering\textit{[Site map placeholder]}
    \end{column}
  \end{columns}
\end{frame}

% ------------------------------------------------------------------
\begin{frame}{Pipeline Architecture}
  \begin{enumerate}
    \item \textbf{L1 --- Raw to QC}\\
          Ingest burst files $\rightarrow$ clock drift correction
          $\rightarrow$ quality control $\rightarrow$ tilt correction
          $\rightarrow$ coordinate rotation
    \item \textbf{L2 --- Spectral Analysis}\\
          Segmentation $\rightarrow$ Welch PSD / cross-spectra
          $\rightarrow$ pressure correction $\rightarrow$ bulk
          parameters, bed velocity, stresses
    \item \textbf{Validation}\\
          PUV vs.\ CDIP MOP wave model, spectral bias investigation
    \item \textbf{Analysis scripts}\\
          Cross-deployment diagnostics, hypothesis testing
  \end{enumerate}
  \vspace{0.5em}
  \footnotesize
  GitHub: \texttt{holdenlesliebole/Nortek\_Vector\_PUV\_Pipeline}
\end{frame}

% ==================================================================
\section{L1 Processing}
% ==================================================================

% ------------------------------------------------------------------
\begin{frame}{Raw Data Ingestion}
  \begin{itemize}
    \item Nortek Vector burst files: \texttt{.dat} (data),
          \texttt{.sen} (sensor/diagnostics), \texttt{.hdr} (header)
    \item Parsed with MATLAB \texttt{textscan} for speed ---
          deployments can exceed $10^8$ samples
    \item Continuous \SI{2}{\hertz} time series reconstructed from
          burst boundaries
    \item \textbf{Clock drift correction}: linear interpolation
          between deployment and recovery clocks
    \item Battery cutoff detection: gaps $> \SI{1}{\second}$
          $\Rightarrow$ discard remainder of record
  \end{itemize}
\end{frame}

% ------------------------------------------------------------------
\begin{frame}{Quality Control --- Old vs.\ New Approach}
  \begin{columns}[T]
    \begin{column}{0.48\textwidth}
      \textbf{Old: Fixed threshold}
      \begin{itemize}
        \item Absolute $\pm\SI{5}{\degree}$ pitch/roll cap
        \item Discards all tilted data
        \item Problematic when pipes bend during storms ---
              stable but non-zero tilt is \emph{correctable}
      \end{itemize}
    \end{column}
    \begin{column}{0.48\textwidth}
      \textbf{New: Variability-based QC}
      \begin{itemize}
        \item Rolling std of pitch/roll over \SI{60}{\second} windows
        \item Flag samples where variability $> \SI{2}{\degree}$
        \item Absolute cap retained at $\SI{30}{\degree}$ for
              severely compromised instruments
        \item Stable tilts \emph{retained} and corrected
      \end{itemize}
    \end{column}
  \end{columns}
\end{frame}

% ------------------------------------------------------------------
\begin{frame}{Tilt Correction}
  Sample-by-sample 3D rotation using measured pitch ($\phi$) and
  roll ($\theta$):
  %
  \begin{equation*}
    \mathbf{R} = \mathbf{R}_\text{roll}(\theta)\;\mathbf{R}_\text{pitch}(\phi)
    = \begin{bmatrix}
        1 & 0 & 0 \\
        0 & \cos\theta & -\sin\theta \\
        0 & \sin\theta &  \cos\theta
      \end{bmatrix}
      \begin{bmatrix}
        \cos\phi & 0 & \sin\phi \\
        0 & 1 & 0 \\
        -\sin\phi & 0 & \cos\phi
      \end{bmatrix}
  \end{equation*}
  %
  \begin{itemize}
    \item Applied \emph{before} heading rotation to buoy coordinates
    \item Corrects velocity components for instrument tilt
    \item Critical for instruments with bent mounting pipes
          (post-storm)
  \end{itemize}
\end{frame}

% ------------------------------------------------------------------
\begin{frame}{Coordinate Rotation}
  \begin{itemize}
    \item Instrument XYZ $\rightarrow$ geographic buoy coordinates
    \item Convention: $+x$ West, $+y$ North, $+z$ Up (left-handed,
          MOP-consistent)
    \item Heading from deployment notes $+$ magnetic declination via
          IGRF-13 model at deployment location and epoch
    \item Shore-normal rotation applied later in L2 using CDIP
          THREDDS shore-normal angle for each MOP station
  \end{itemize}
\end{frame}

% ------------------------------------------------------------------
\begin{frame}{L1 Processing Results}
  \begin{center}
    \Large \textbf{33 / 40} instruments processed successfully
    (\textbf{82.5\%})
  \end{center}
  \vspace{0.5em}
  \textbf{7 failures} --- all hardware-related:
  \begin{itemize}
    \item Instrument burial (sand accretion)
    \item Bent or broken mounting pipes
    \item Dead / depleted batteries (no usable data)
  \end{itemize}
  \vspace{0.5em}
  No failures attributable to the processing pipeline itself.
\end{frame}

% ==================================================================
\section{L2 Spectral Analysis}
% ==================================================================

% ------------------------------------------------------------------
\begin{frame}{Segmentation}
  \begin{itemize}
    \item \textbf{Segment length}: 2048 samples = \SI{17.1}{\minute}
          at \SI{2}{\hertz}
    \item Non-overlapping segments
    \item \textbf{NaN policy}: segments with $>$10\% NaN discarded;
          $\leq$10\% filled by linear interpolation
    \item Segment-mean pressure $\rightarrow$ water depth $H$
    \item All channels linearly detrended before spectral analysis
  \end{itemize}
\end{frame}

% ------------------------------------------------------------------
\begin{frame}{Spectral Estimation}
  \begin{itemize}
    \item \textbf{Method}: Welch's averaged periodogram
    \item Sub-segments: $N_\text{fft} = 256$ samples (\SI{128}{\second})
    \item \textbf{Window}: Hanning, 50\% overlap
    \item $\approx$\textbf{34 effective DOF} per segment
    \item Frequency resolution: $\Delta f \approx
          \SI{0.0078}{\hertz}$, range 0--\SI{1}{\hertz}
    \item \textbf{Cross-spectra}: $S_{pu}$, $S_{pv}$ computed with
          same windowing for directional coefficients $a_1$, $b_1$
  \end{itemize}
\end{frame}

% ------------------------------------------------------------------
\begin{frame}{Pressure Correction}
  Surface elevation spectrum from pressure:
  %
  \begin{equation*}
    S_\eta(f) = \frac{S_p(f)}{K_p^2(f)}, \qquad
    K_p(f) = \frac{\cosh(k\, z_s)}{\cosh(k\, H)}
  \end{equation*}
  %
  \begin{itemize}
    \item Wavenumber $k(f)$: linear dispersion relation
          $\omega^2 = g k \tanh(kH)$ solved via Newton's method
    \item \textbf{Cutoff}: $K_p < 0.1$ $\Rightarrow$ set
          $S_\eta(f) = 0$ (noise suppression)
    \item Resulting $f_\text{cut}$ recorded per segment
    \item Verified: Newton, Wu (1986), and O'Reilly formulations
          produce identical $K_p$ (see Validation section)
  \end{itemize}
\end{frame}

% ------------------------------------------------------------------
\begin{frame}{Bulk Wave Parameters}
  \begin{itemize}
    \item $H_s = 4\sqrt{m_0}$, computed separately for:
      \begin{itemize}
        \item Sea--swell: \SIrange{0.04}{0.25}{\hertz}
        \item Infragravity: \SIrange{0.004}{0.04}{\hertz}
      \end{itemize}
    \item Peak period $T_p = 1/f_p$ (max $S_\eta$ in SS band)
    \item Mean period $T_{m02} = \sqrt{m_0/m_2}$
    \item Mean direction from $a_1$, $b_1$ (PUV cross-spectra)
    \item Energy flux
          $F = \rho g \int c_g(f)\, S_\eta(f)\, df$
  \end{itemize}
\end{frame}

% ------------------------------------------------------------------
\begin{frame}{Additional L2 Products}
  \begin{itemize}
    \item \textbf{Near-bed velocity}: IFFT method ---
          scale velocity FFT by
          $\cosh(kH)/\cosh[k(H - z_s)]$, inverse transform
    \item \textbf{Bed shear stress}: Swart (1974) friction factor,
          $k_s = 2.5\, D_{50}$
    \item \textbf{Reynolds stress}: $\overline{u'w'}$,
          $\overline{v'w'}$, $\overline{u'v'}$; turbulent KE
    \item \textbf{Velocity moments}: skewness $S_k$, asymmetry $A_s$
          (Elgar 1985; Hoefel \& Elgar 2003)
  \end{itemize}
\end{frame}

% ------------------------------------------------------------------
\begin{frame}{Shore-Normal Rotation}
  \begin{itemize}
    \item Shore-normal angle retrieved from CDIP THREDDS for each
          MOP station
    \item Velocities rotated: $+x$ onshore, $+y$ alongshore north
    \item \textbf{Fallback}: if THREDDS unavailable, retain
          geographic buoy coordinates ($+x$ W, $+y$ N)
    \item Shore-normal frame required for cross-shore energy flux
          and transport calculations
  \end{itemize}
\end{frame}

% ==================================================================
\section{Validation Against MOP}
% ==================================================================

% ------------------------------------------------------------------
\begin{frame}{MOP Wave Model}
  \begin{itemize}
    \item CDIP Monitoring and Prediction (MOP) spectral wave model
    \item Initialized by offshore CDIP buoy observations
    \item Provides $S_\eta(f)$ at \SI{10}{\meter} depth, hourly
    \item \textbf{Shoaling to PUV depth}:
  \end{itemize}
  \begin{equation*}
    S_\eta(f,\, h_\text{PUV}) = S_\eta(f,\, 10\,\text{m})
    \times \frac{c_g(f,\, 10\,\text{m})}{c_g(f,\, h_\text{PUV})}
  \end{equation*}
  \begin{itemize}
    \item PUV segments (\SI{17}{\minute}) interpolated to MOP hourly
          time base for comparison
  \end{itemize}
\end{frame}

% ------------------------------------------------------------------
\begin{frame}{Bulk Comparison Results}
  \begin{columns}[T]
    \begin{column}{0.50\textwidth}
      \begin{itemize}
        \item $H_s$: $R^2 = 0.83$--$0.86$
        \item RMSE $< \SI{0.06}{\meter}$
        \item Small positive bias (PUV $>$ MOP)
        \item Kelp-fouled instrument: $R^2 = 0.66$,
              bias $= \SI{+0.037}{\meter}$
      \end{itemize}
    \end{column}
    \begin{column}{0.48\textwidth}
      \includegraphics[width=\textwidth]{TBR23_MOP580_7m_PUV_vs_MOP.png}
    \end{column}
  \end{columns}
\end{frame}

% ------------------------------------------------------------------
\begin{frame}{Frequency-Dependent Bias}
  \begin{columns}[T]
    \begin{column}{0.50\textwidth}
      \begin{itemize}
        \item Excess energy concentrated at \textbf{swell peak}
              (\SIrange{0.04}{0.08}{\hertz})
        \item $K_p > 0.9$ at these frequencies --- pressure
              correction is minimal
        \item Mid-swell band: near-unity ratio (0.99--1.03)
        \item \textbf{NOT} a high-frequency pressure correction
              artifact
      \end{itemize}
    \end{column}
    \begin{column}{0.48\textwidth}
      \includegraphics[width=\textwidth]{TBR23_spectral_ratio_all_instruments.png}
      % fallback: TBR23_band_ratio_vs_depth.png
    \end{column}
  \end{columns}
\end{frame}

% ------------------------------------------------------------------
\begin{frame}{Pressure Correction Investigation}
  \begin{columns}[T]
    \begin{column}{0.50\textwidth}
      \begin{itemize}
        \item Compared three wavenumber formulations:
          \begin{enumerate}
            \item Newton's method (baseline)
            \item Wu (1986) approximation
            \item O'Reilly coefficients
          \end{enumerate}
        \item All produce \textbf{identical} $K_p$
        \item Pressure correction is \textbf{not} the source
              of the observed bias
      \end{itemize}
    \end{column}
    \begin{column}{0.48\textwidth}
      \includegraphics[width=\textwidth]{wavenumber_comparison_H6m.png}
    \end{column}
  \end{columns}
\end{frame}

% ==================================================================
\section{Hypothesis Testing}
% ==================================================================

% ------------------------------------------------------------------
\begin{frame}{Overview: Four Hypotheses}
  What causes the systematic positive $H_s$ bias (PUV $>$ MOP)?
  \vspace{0.5em}
  \begin{enumerate}
    \item \textbf{Nonlinear shoaling} --- triad interactions
          (Ursell number)
    \item \textbf{Directional narrowing} --- refraction-driven
          focusing
    \item \textbf{Bound long waves} --- group-forced IG
          contamination
    \item \textbf{MOP spectral peak broadening} --- numerical
          diffusion in the wave model
  \end{enumerate}
\end{frame}

% ------------------------------------------------------------------
\begin{frame}{H1: Nonlinear Shoaling}
  \begin{columns}[T]
    \begin{column}{0.50\textwidth}
      \begin{itemize}
        \item Test: correlate spectral ratio with Ursell number
              $\mathit{Ur} = H_s L_p^2 / h^3$
        \item $N = 5228$ segment-hours
        \item \textbf{Result}: $R^2 = 0.004$
        \item Median $\mathit{Ur}$: 12 at \SI{8.3}{\meter},
              30 at \SI{5.5}{\meter}
        \item[$\Rightarrow$] \alert{RULED OUT}
      \end{itemize}
    \end{column}
    \begin{column}{0.48\textwidth}
      \includegraphics[width=\textwidth]{TBR23_ursell_vs_ratio.png}
    \end{column}
  \end{columns}
\end{frame}

% ------------------------------------------------------------------
\begin{frame}{H2: Directional Narrowing}
  \begin{itemize}
    \item Refraction should \emph{narrow} directional spread during
          shoaling $\Rightarrow$ energy focusing $\Rightarrow$ higher
          $H_s$
    \item \textbf{Result}: PUV spread is systematically
          \emph{wider} than MOP at swell frequencies
          (\SI{15}{\degree} vs.\ \SI{10}{\degree})
    \item Opposite to prediction
    \item Correlation with $H_s$ amplification: $R^2 < 0.001$
    \item[$\Rightarrow$] \alert{RULED OUT}
  \end{itemize}
  \vspace{0.5em}
  \footnotesize Wider PUV spread may reflect scattered energy at
  shallow depth not captured by the forward-propagating MOP model.
\end{frame}

% ------------------------------------------------------------------
\begin{frame}{H3: Bound Long Waves}
  \begin{columns}[T]
    \begin{column}{0.50\textwidth}
      \begin{itemize}
        \item Bound wave energy $\propto H_s^2$
              (Longuet-Higgins \& Stewart 1962)
        \item \textbf{Result}: no positive correlation with
              $H_s^2$ ($R = -0.14$)
        \item Excess in \emph{wrong frequency band} ---
              at swell peak (0.04--0.08~Hz), not at group
              frequencies (0.01--0.03~Hz)
        \item[$\Rightarrow$] \alert{RULED OUT}
      \end{itemize}
    \end{column}
    \begin{column}{0.48\textwidth}
      \includegraphics[width=\textwidth]{TBR23_MOP580_7m_bound_wave_scaling.png}
    \end{column}
  \end{columns}
\end{frame}

% ------------------------------------------------------------------
\begin{frame}{H4: MOP Spectral Peak Broadening}
  \begin{columns}[T]
    \begin{column}{0.50\textwidth}
      \begin{itemize}
        \item $Q_p$ (Goda peakedness): PUV 7\% higher than MOP
        \item Half-power bandwidth: PUV \textbf{16.7\% narrower}
        \item Peak spectral density: PUV 14.4\% higher
        \item $\Delta Q_p$ vs.\ $H_s$ amplification:
              $R = 0.44$, $R^2 = 0.20$
        \item[$\Rightarrow$] \textbf{\color{mLightGreen}SUPPORTED}
      \end{itemize}
    \end{column}
    \begin{column}{0.48\textwidth}
      \includegraphics[width=\textwidth]{TBR23_MOP580_7m_normalized_shape.png}
    \end{column}
  \end{columns}
  \vspace{0.3em}
  \footnotesize MOP model introduces numerical broadening of the
  spectral peak during propagation; PUV captures the true, narrower
  peak shape.
\end{frame}

% ------------------------------------------------------------------
\begin{frame}{Peakedness Diagnostic}
  \begin{center}
    \includegraphics[width=0.65\textwidth]{TBR23_MOP580_7m_peakedness.png}
  \end{center}
  Hours when PUV spectrum is more peaked relative to MOP correspond
  to hours with larger positive $H_s$ bias.
\end{frame}

% ------------------------------------------------------------------
\begin{frame}{Hypothesis Summary}
  \begin{table}
    \centering
    \small
    \begin{tabular}{@{}llll@{}}
      \toprule
      \textbf{Hypothesis} & \textbf{Test metric} &
        \textbf{Result} & \textbf{Verdict} \\
      \midrule
      1. Nonlinear shoaling &
        $R^2(\mathit{Ur},\, \text{ratio})$ &
        $R^2 = 0.004$ & Ruled out \\
      2. Directional narrowing &
        PUV vs.\ MOP spread &
        PUV wider & Ruled out \\
      3. Bound long waves &
        $R(H_s^2,\, \text{excess})$ &
        $R = -0.14$ & Ruled out \\
      4. Peak broadening &
        $R(\Delta Q_p,\, \text{bias})$ &
        $R = 0.44$ & \textbf{Supported} \\
      \bottomrule
    \end{tabular}
  \end{table}
\end{frame}

% ------------------------------------------------------------------
\begin{frame}{Cross-Deployment Confirmation}
  \begin{itemize}
    \item Pattern confirmed across \textbf{11 instruments}
          from multiple deployments
    \item Correlation $R(Q_p) = 0.30$--$0.63$ across instruments
    \item Spans:
      \begin{itemize}
        \item Summer and winter wave climates
        \item Depths \SIrange{5.5}{15.5}{\meter}
        \item Multiple MOP transects
      \end{itemize}
    \item Spectral peak broadening is a \textbf{robust, general}
          feature of the MOP model --- not site- or season-specific
  \end{itemize}
\end{frame}

% ==================================================================
\section{Key Findings and Next Steps}
% ==================================================================

% ------------------------------------------------------------------
\begin{frame}{Key Findings}
  \begin{enumerate}
    \item \textbf{Pipeline validated}: 33/40 instruments processed
          (82.5\%), strong agreement with MOP
          ($R^2 > 0.83$, RMSE $< \SI{0.06}{\meter}$)
    \item \textbf{Bias identified}: small systematic positive
          $H_s$ bias concentrated at the swell spectral peak,
          \emph{not} from the pressure correction
    \item \textbf{Mechanism determined}: MOP model introduces
          numerical spectral peak broadening; PUV captures the true
          narrower peak ($R = 0.44$ between $\Delta Q_p$ and bias)
    \item \textbf{Cross-deployment robustness}: pattern confirmed
          across 11 instruments, summer/winter, 5.5--15.5~m depth
  \end{enumerate}
\end{frame}

% ------------------------------------------------------------------
\begin{frame}{Implications for TBR23 Paper}
  \begin{itemize}
    \item Pipeline is validated and ready for production use
    \item Bias is understood and attributable to the model, not the
          measurement --- no correction needed for PUV data
    \item Spectral shape finding \emph{supports} using PUV as
          ground truth for nearshore wave characterization
    \item L2 products (bed velocity, stress, moments) are ready
          for sediment transport analysis
  \end{itemize}
\end{frame}

% ------------------------------------------------------------------
\begin{frame}{Next Steps}
  \begin{itemize}
    \item \textbf{Cross-deployment analysis}: expand when CDIP
          THREDDS recovers (currently intermittent)
    \item \textbf{IG / bound-wave quantification}: separate forced
          vs.\ free IG energy
    \item \textbf{Tidal decomposition}: segment analysis by tidal
          phase and water level
    \item \textbf{Grain size integration}: laser diffraction
          $D_{50}$ for site-specific bed roughness
    \item \textbf{L3 transport calculations}: combine wave
          statistics with morphology data
    \item \textbf{Potential spinoff paper}: wave spectral shape
          evolution during shoaling --- novel PUV-based
          quantification of model spectral broadening
  \end{itemize}
\end{frame}

% ------------------------------------------------------------------
\begin{frame}[standout]
  Questions?
\end{frame}

\end{document}
